% !TeX root = ../thesis.tex
\chapter{Conclusion}
\label{chap:conclusion}

This thesis built and evaluated a weakly supervised pipeline that answers three
questions for each UCF-Crime video: whether it is anomalous, which of three
project-defined services it involves, and where the event occurs. The pipeline
adapts VideoMAE with LoRA, adds a normal-prototype distance to frozen clip
features, and trains an attention MIL model with a binary head and a
three-channel response head. Each output was compared with a trivial rule on the
same videos, across repeated Stage-2 training and two test input paths.

The three outputs behave very differently. Video-level anomaly detection is
strong and stable across seeds and input paths. Response selection is weak: the
response head defaults to Police and over-predicts the other services, and after
the binary gate it adds little over predicting every service for each detected
anomaly. Adding normal bags to the response loss removes the Police default but
not the weak selection. Attention ranks annotated frames above trivial
references, but its intervals are no more precise than predicting the whole
video.

The contribution is therefore less a better detector than a clearer measurement
of what such a pipeline delivers. Matched trivial rules, per-video temporal
references, and the separation of seed variation from test-sampling uncertainty
showed that most of the apparent response capability comes from the binary
decision. Response targets derived from dataset categories reward a default
rather than selective behavior; meaningful response prediction needs targets
grounded in expert judgment or observed service needs.

The next steps follow directly: expert-validated response targets with an
abstention option, matched component baselines on the existing pipeline,
evaluation on external data, and grounding studied as its own task. Operational
integration, with human oversight and privacy governance, belongs only after
these steps.
